\section{System Taxonomy: Agents and Architecture}

Before detailing the computational processes of Collective, we must define the primary entities that inhabit and sustain the mesh. Unlike the rigid hierarchies of the legacy market, roles in Collective are fluid, determined by current activity rather than permanent employment status.

\subsection{Community Members (The Workers)}
The heart of the Collective is the \textbf{Community Member}, referred to functionally as the \textbf{Worker} during task execution. Members are the sole "processors" of the system. 
\begin{itemize}[leftmargin=*]
    \item \textbf{Sovereignty:} Members own their data and move freely between different Service Nodes and geographic instances.
    \item \textbf{Identity:} A Member's identity is tied to their global "Body of Knowledge" certifications and their accumulated Value balance.
\end{itemize}

\subsection{Service Nodes and Specialized Pillars}
The architecture is composed of recursive, specialized entities that provide the "Operating System" services:
\begin{enumerate}
    \item \textbf{Service Nodes:} The primary units of production (e.g., an Eco-Village bakery).
    \item \textbf{Bodies of Knowledge (BoK):} The educational kernel; they define what a Task is and verify if a Member possesses the required Skill.
    \item \textbf{Societies:} The legal and territorial layer; they govern the physical "Lease" of the land and protect member rights.
    \item \textbf{The Common Fund:} The metabolic gateway; it manages external procurement and internal liquidity.
\end{enumerate}

\subsection{Active Agents: Designers and Prospectors}
While every member is a Worker, some choose to act as \textbf{Service Designers} (the architects of the WBS) or \textbf{Demand Prospectors} (the sensory scouts identifying community needs). These roles allow the mesh to proactively evolve its production logic.